\documentclass[12pt]{article}
\usepackage[margin=1in]{geometry}
\usepackage{times}
\usepackage{caption}
\usepackage{booktabs}
\usepackage{array}
\usepackage{tikz}
\usetikzlibrary{positioning, shapes.geometric, arrows}

\begin{document}

Geele Evans\\
Dr. Lori Jirousek-Falls\\
FCWR 304-F02\\
25 March 2026

Remote Work Security: Protecting Home Networks for Small Business Employees

Abstract
Remote work expands opportunity for small businesses but shifts risk to home networks. This memo translates current CISA and NIST guidance into a concise, actionable framework for managers with no dedicated IT staff. It identifies router vulnerabilities, the importance of VPNs, Wi‑Fi hardening, and day‑to‑day practices, then provides a practical implementation checklist that can be deployed this week. The goal is to establish a secure baseline that is auditable, scalable, and affordable.

Threat Landscape and Rationales
Small businesses face router vulnerabilities from outdated firmware, weak credentials, and misconfigurations; VPN exposure when not used or poorly configured; and Wi‑Fi protections that are insufficient or improperly segmented. When home networks are compromised, work devices and corporate data are at risk. A practical program for small businesses should emphasize tangible improvements, not theory. This framing aligns with CISA and NIST guidance that emphasizes perimeter hardening and verifiable controls (CISA; NIST).

Practical Fixes: Technical and Policy Measures

Router Security
- Keep firmware current; enable automatic updates if possible.
- Change default admin credentials to a strong, unique password.
- Disable remote management unless necessary; restrict if used.
- Turn off UPnP and WPS; enable the router firewall.
- Create a separate work network (dedicated SSID) for company devices.

VPN and Remote Access
- Mandate VPN use for all remote work; require MFA for VPN access.
- Establish clear VPN usage policies, including device posture checks before connection.
- Minimize or eliminate split tunneling to reduce exposure from home devices (CISA; NIST guidance).

Wi‑Fi Hardening
- Use WPA3 on the home network; disable older protocols (WEP, WPA2‑TKIP).
- Avoid auto‑connecting to unknown networks; segment work devices on a dedicated network.
- Maintain network segmentation so work traffic stays isolated from personal devices.

Day-to-Day Employee Practices
- Keep devices updated; enable automatic updates.
- Use a password manager and MFA for all work accounts.
- Enable device encryption and a short automatic screen lock timeout.
- If using public networks, employ a VPN; avoid untrusted networks for work tasks.
- Regularly back up work data to approved storage with versioning.

Implementation Checklist for Employees
- Update router firmware; change admin password; disable remote management; disable UPnP/WPS; create a dedicated work network.
- Enforce WPA3 on the work network with a strong passphrase.
- Install and configure a company‑approved VPN with MFA; route work traffic through the VPN.
- Require pre‑VPN device posture checks (up-to-date OS, active antivirus).
- Enable device encryption and a short auto‑lock timeout.
- Use a password manager and enable MFA on work services.
- Avoid public Wi‑Fi for work; use VPN if unavoidable.
- Back up work data to approved repositories/storage with versioning.
- Report any device compromise or router anomaly immediately.

Figure 1. Threat Landscape Map for Remote Work Security
Caption: A schematic of the principal risk vectors for home networks used by remote workers—router vulnerabilities, VPN exposure, Wi‑Fi misconfigurations, and day-to-day endpoint practices. Adapted from CISA and NIST.
\\begin{figure}[h]
\\centering
\\begin{tikzpicture}[>=stealth', node distance=3cm, auto]
\\node[draw, rounded corners, minimum width=3.6cm, minimum height=1.6cm] (r) {Router Vulnerabilities};
\\node[draw, rounded corners, right=of r, minimum width=3.6cm, minimum height=1.6cm] (w) {Work Devices};
\\node[draw, rounded corners, below=of r, minimum width=3.6cm, minimum height=1.6cm] (v) {VPN Exposure};
\\node[draw, rounded corners, right=of v, minimum width=3.6cm, minimum height=1.6cm] (c) {Corporate Data};
\\node[draw, rounded corners, left=of r, minimum width=3.6cm, minimum height=1.6cm] (b) {WiFi Misconfig};
\\node[draw, rounded corners, below=of b, minimum width=3.6cm, minimum height=1.6cm] (e) {End Points};

\\draw[->] (r) -- (w);
\\draw[->] (v) -- (c);
\\draw[->] (b) -- (e);
\\draw[->] (e) -- (c);
\\end{tikzpicture}
\\caption{Threat Landscape Map for Remote Work Security.}
\\end{figure}

Figure 2. VPN Topology Diagram
Caption: A recommended remote-work topology: home device connects to a manufacturer‑validated VPN with MFA, routing all work traffic through the corporate network.
\\begin{figure}[h]
\\centering
\\begin{tikzpicture}[>=stealth', node distance=3cm, auto]
\\node[ellipse, draw] (home) {Home Device};
\\node[ellipse, draw, right=of home] (vpn) {VPN Gateway};
\\node[ellipse, draw, right=of vpn] (corp) {Corporate Network};

\\path[->] (home) edge (vpn);
\\path[->] (vpn) edge (corp);
\\path[->] (home) edge[bend left=25] (corp) node[midway, above] {Internet};
\\end{tikzpicture}
\\caption{VPN Topology Diagram.}
\\end{figure}

Table 1. Day-to-Day Employee Practices for Remote Security
Caption: Practical daily practices aligned with CISA/NIST guidance.
\\begin{table}[h]
\\centering
\\begin{tabular}{p{4.2cm} p{5.5cm} p{2.2cm} p{2.2cm} p{2.4cm}}
\\toprule
Practice & Rationale & Owner & Frequency & Status \\
\\midrule
Update OS & Closes known vulnerabilities & Employee & Auto/as scheduled & Pending \\
Use MFA on all work accounts & Reduces credential abuse & Employee & Per account & In progress \\
Use a password manager & Encourages strong passwords & Employee & Always & Required \\
Enable device encryption & Protects data at rest & Employee & On devices & Recommended \\
Require VPN for work & Encapsulates traffic; lowers exposure & All users & Always & Mandatory \\
Avoid public Wi‑Fi for work & Reduces risk from untrusted networks & Employee & Always & Mandatory \\
Back up work data regularly & Preserves data integrity & Employee/Manager & Daily/weekly & Ongoing \\
\\bottomrule
\\end{tabular}
\\end{table}

Works Cited
CISA. “Remote Work Security Guidance.” Cybersecurity and Infrastructure Security Agency, 2023. Accessed 25 Mar. 2026.
NIST. “Guidelines for Securing Home Networks.” National Institute of Standards and Technology, 2020. Accessed 25 Mar. 2026.
Evans, Geele. “Remote Work Security: Protecting Home Networks for Small Business Employees.” FCWR 304 Memo, 6 Mar. 2026.